% !TEX root = ../../hw03.tex
% Homework 03, question 4.
\begin{proof}[Solution]
  For part (a), the initial values and the recurrence show by induction
  that $F_n\geq0$ for every $n\geq0$, and $F_n\geq1$ for every
  $n\geq1$. For $n\geq1$, the recurrence gives
  \[
    F_{n+1}-F_n=F_{n-1}\geq0.
  \]
  Together with $F_1\geq F_0$, this proves that the sequence is
  nondecreasing. It is strictly increasing from $F_2$ onward, since
  $F_{n-1}>0$ when $n\geq2$. The only repeated adjacent values are
  $F_1=F_2=1$.

  % This proof could be much clearer
  % For part (b), we have $F_{n+1}\geq F_n+1$ whenever $n\geq2$.
  % Starting from $F_2=1$, induction therefore gives
  % \[
  %   \forall n\geq2,\quad F_n\geq n-1.
  % \]
  % Given any real number $M$, the Archimedean property supplies an
  % integer $n\geq2$ with $n-1>M$. Then $F_n>M$, so the sequence is
  % unbounded above.

  For part (b), suppose $(F_n)$ is bounded above. Since it is nondecreasing,
  the monotone convergence theorem gives $F_n\to L$ as $n\to\infty$.
  The shifted sequences $F_{n-1}$ and $F_{n-2}$ also tend to $L$, so the
  recurrence and the limit laws give
  \begin{align*}
    F_n&=F_{n-1}+F_{n-2} \\
       &\implies L=2L \\
       &\iff L=0.
  \end{align*}
  But $F_n\geq1$ for every $n\geq1$, so taking limits gives $L\geq1$.
  This contradicts $L=0$. Therefore the sequence is unbounded above.

  For part (c), the ratios are well defined because their denominators
  are positive. We have $r_2=1$. For $n\geq3$, substitute the recurrence
  into the numerator and split the fraction to obtain
  \begin{align*}
    r_n&=\frac{F_n}{F_{n-1}}\\
       &=\frac{F_{n-1}+F_{n-2}}{F_{n-1}}\\
       &=\frac{F_{n-1}}{F_{n-1}}+\frac{F_{n-2}}{F_{n-1}}\\
       &=1+\frac{F_{n-2}}{F_{n-1}}.
  \end{align*}
  By part (a), we have $0<F_{n-2}\leq F_{n-1}$. Dividing by the
  positive number $F_{n-1}$ gives
  \[
    0<\frac{F_{n-2}}{F_{n-1}}\leq1.
  \]
  Adding $1$ now gives $1<r_n\leq2$.
  Including $r_2$, we conclude that
  \[
    \forall n\geq2,\quad 1\leq r_n\leq2.
  \]
  Thus the ratio sequence is bounded.

  For part (d), assume that $r_n\to L$. The bounds above give
  $1\leq L\leq2$, so $L$ is nonzero.
  \begin{marginfigure}
    \centering
    % Equal axis units preserve the geometry of the diagonal.
% Curve labels sit in a separate legend so the intersection has clear space.
\begin{tikzpicture}[interval diagram,x=12mm,y=12mm]
  \pgfmathsetmacro{\goldenratio}{(1+sqrt(5))/2}

  \draw[->] (0,0) -- (2.8,0)
    node[right=4pt,interval label] {$x$};
  \draw[->] (0,0) -- (0,3.3)
    node[above=4pt,interval label] {$y$};
  \node[below left=4pt,interval label] at (0,0) {$0$};

  \draw[gray!65,densely dotted]
    (\goldenratio,0) -- (\goldenratio,\goldenratio)
    -- (0,\goldenratio);
  \draw (\goldenratio,-2pt) -- (\goldenratio,2pt);
  \draw (-2pt,\goldenratio) -- (2pt,\goldenratio);
  \node[below=5pt,interval label] at (\goldenratio,0) {$\varphi$};
  \node[left=5pt,interval label] at (0,\goldenratio) {$\varphi$};

  \draw[gray!55,dashed,->] (0,0) -- (2.55,2.55);
  \draw[<->,domain=0.46:2.55,samples=80]
    plot (\x,{1+1/\x});
  \fill (\goldenratio,\goldenratio) circle[radius=1.6pt];

  \draw[gray!55,dashed] (0,-0.75) -- (0.45,-0.75);
  \node[anchor=west,interval label] at (0.65,-0.75) {$y=x$};
  \draw (0,-1.3) -- (0.45,-1.3);
  \node[anchor=west,interval label] at (0.65,-1.3) {$y=1+1/x$};
\end{tikzpicture}

    \caption[The golden ratio as a fixed point.]{The positive intersection
      satisfies $\varphi=1+1/\varphi$. Dotted guides locate both coordinates
      at $\varphi\approx1.618$. The axes use equal units.}
    \label{fig:hw03:golden-ratio-fixed-point}
  \end{marginfigure}
  To express $r_n$ in terms of $r_{n-1}$, recall from part (c) that,
  for $n\geq3$,
  \[
    r_n=\frac{F_n}{F_{n-1}}=1+\frac{F_{n-2}}{F_{n-1}}.
  \]
  The preceding ratio is $r_{n-1}=F_{n-1}/F_{n-2}$. Since both
  Fibonacci terms are positive, we can take its reciprocal to obtain
  \[
    \frac{1}{r_{n-1}}
      =\frac{1}{F_{n-1}/F_{n-2}}
      =\frac{F_{n-2}}{F_{n-1}}.
  \]
  Substituting this into the expression for $r_n$ gives the recurrence
  \[
    r_n=1+\frac{1}{r_{n-1}},\qquad n\geq3.
  \]
  The shifted sequence $r_{n-1}$ also tends to $L$. Since $L\neq0$,
  the limit laws give
  \[
    L=1+\frac{1}{L}.
  \]
  Geometrically, $(L,L)$ lies on both $y=x$ and $y=1+1/x$, as shown
  in \autoref{fig:hw03:golden-ratio-fixed-point}.
  Multiplying by $L$ and bringing all terms to one side gives
  \[
    L^2=L+1,\qquad L^2-L-1=0.
  \]
  The two roots are $(1+\sqrt{5})/2$ and $(1-\sqrt{5})/2$.
  The second is negative and cannot satisfy $L\geq1$. Therefore the
  assumed limit is the golden ratio
  \[
    L=\frac{1+\sqrt{5}}{2}=\varphi\approx1.6180339887.
  \]
\end{proof}

The identity $\varphi-1=1/\varphi$ also has a geometric interpretation.
\begin{marginfigure}
  \centering
  % A golden rectangle and its recursive square subdivision.
% The black curve is a logarithmic golden spiral, not circular arcs.
% Geometry reference: https://doi.org/10.1155/2019/3149602
\begin{tikzpicture}[interval diagram,x=28mm,y=28mm]
  \pgfmathsetmacro{\goldenratio}{(1+sqrt(5))/2}
  \pgfmathsetmacro{\shrink}{1/\goldenratio}
  \pgfmathsetmacro{\eyex}{\goldenratio/(1+\shrink^2)}
  \pgfmathsetmacro{\eyey}{\shrink^2/(1+\shrink^2)}

  % A clockwise quarter-turn and contraction reproduce the next square.
  \foreach \step in {0,...,5} {
    \pgfmathsetmacro{\factor}{\shrink^\step}
    \pgfmathsetmacro{\turn}{-90*\step}
    \begin{scope}[
      shift={(\eyex,\eyey)},rotate=\turn,scale=\factor,
      shift={(-\eyex,-\eyey)}
    ]
      \draw[black!35] (0,0) rectangle (1,1);
    \end{scope}
  }
  \draw[black!55] (0,0) rectangle (\goldenratio,1);

  % Contract continuously while turning clockwise around the same centre.
  \draw[domain=0:630,samples=220,variable=\angle]
    plot (
      {\eyex+\shrink^(\angle/90)*
        (-\eyex*cos(\angle)-\eyey*sin(\angle))},
      {\eyey+\shrink^(\angle/90)*
        (\eyex*sin(\angle)-\eyey*cos(\angle))}
    );

  % Only label the two largest squares; leave the inner turns uncluttered.
  \node[interval label] at (0.43,0.39) {$1$};
  \node[interval label] at (1.3,0.68) {$\varphi^{-1}$};
  \node[above=6pt,interval label] at (\goldenratio/2,1) {$\varphi$};
  \node[left=6pt,interval label] at (0,0.5) {$1$};
\end{tikzpicture}

  \caption[The golden rectangle and golden spiral.]{A golden rectangle
    subdivides into squares with side lengths
    $1,\varphi^{-1},\varphi^{-2},\ldots$. The golden spiral passes through
    successive corners. Each quarter-turn inward shrinks the distance
    from its centre by $1/\varphi$. The subdivision continues indefinitely.}
  \label{fig:hw03:golden-rectangle-spiral}
\end{marginfigure}
A rectangle with side lengths $1$ and $\varphi$ leaves a smaller,
rotated copy of itself when a unit square is removed. Repeating this
step gives the squares in \autoref{fig:hw03:golden-rectangle-spiral}.
The golden spiral winds through their corners, shrinking its distance
from the centre by $1/\varphi$ with each
quarter-turn~\cite{Duan2019GoldenSpirals}.
